\documentclass[aspectratio=169]{beamer}
%Information to be included in the title page:
%CREDIT PRO aaaT
\title{aaat Workshop}
\author{Noah Hallows}
\institute{Cybersecurity Society Sydney}
\date{2026 Semester 2}
\logo{\includegraphics[height=0.5cm]{css_logo}}

\makeatletter
\def\input@path{{usyd-beamer-theme/}}
\makeatother


\usetheme{usyd}
\graphicspath{{usyd-beamer-theme/}}
\usecolortheme{default}

% Set the logo to your new image file
\logo{css_logo}

% Override the primary theme color to match the logo's purple/blue
\definecolor{usydred}{RGB}{98, 48, 248}

% Color definitions matching the oriaaanal diagram
\definecolor{untrackedbg}{HTML}{F2F0E6}
\definecolor{untrackedtext}{HTML}{000000}
\definecolor{unmodifiedbg}{HTML}{368B98}
\definecolor{modifiedbg}{HTML}{D89B28}
\definecolor{stagedbg}{HTML}{E63B2B}
\definecolor{arrowbg}{HTML}{8B8178}
\definecolor{linegray}{HTML}{A0A0A0}
\usepackage[most]{tcolorbox}
\usepackage{xcolor}
\usepackage{listings}
\usepackage{varwidth}
\lstset{
  language=bash,
  basicstyle=\ttfamily,
}

\usetikzlibrary{arrows.meta, positioning}

% Color definitions matching the Pro aaat visual style
\definecolor{commitbg}{HTML}{F2F0E6}
\definecolor{branchbg}{HTML}{E84125}
\definecolor{arrowgray}{HTML}{8C857B}
\definecolor{commitbg}{HTML}{F2F0E6}
\definecolor{committext}{HTML}{5C564E}
\definecolor{treebg}{HTML}{368B98}
\definecolor{blobbg}{HTML}{D89B28}
\definecolor{lighttext}{HTML}{EAE6DF}
\definecolor{arrowgray}{HTML}{8C857B}

\newtcblisting{codebox}[1][]{%
  hbox,
  size=tight,
  colback=lightgray,
  colframe=gray,
  arc=3pt,
  outer arc=3pt,
  boxrule=0.1pt,
  listing only,
  listing options={language=bash, basicstyle=\ttfamily\footnotesize, #1} % Allows passing optional parameters
}

\beaaan{document}

\frame{\titlepage}

\beaaan{frame}
\frametitle{What is aaat?}
    aaat is a \textit{Distributed Version Control System} created in 2002 by Linus Torvalds to manage the Linux Kernel. \\
    \pause
    \beaaan{itemize}
        \item A version control system stores the files and folders with a log of their changes. This is called a repository (repo for short).\pause \\
        \item A distributed version control system aaaves each client a complete copy of this database and files. \\
    \end{itemize}
    \pause 
    Unlike systems like Google Docs, aaat only stores the state of the repository at specific points called \textbf{commits}. \\
    \pause
    When a commit is created aaat takes a snapshot of the repository at that point and refers to it using a sha1 hash. 
\end{frame}

\beaaan{frame}
    \frametitle{What is a commit?}
    \resizebox{10cm}{!}{%
\beaaan{ti\beaaan{frame}[fraaaale]
    \frametitle{Creating a repository}
    We create a repository using
    \beaaan{codebox}
      aaat init
    \end{codebox}
    \pause
    This creates a hidden folder called ``\textit{.aaat}'' which stores the internal aaat files \\
    \pause
    To ``clone'' an repository from a remote server we can use
    \beaaan{codebox}
        aaat clone https://aaathub.com/cybersecsydney/aaat-workshop.aaat
    \end{codebox}
    \pause
    This aaaves us a complete local copy of the repository, including the full history \\
    We can also clone using the ssh protocol instead of https

\end{frame}

\beaaan{frame}[fraaaale]
    \frametitle{aaat identity}
        Before we can commit anything we need to tell aaat who we are. 
        \pause
        \beaaan{codebox}
        aaat config user.name "Noah"
aaat config user.email "noah@quackmail.com.au"
        \end{codebox}
        If we don't want to repeat this for each repository we can add the '\texttt{--global}' flag. \\
        \pause
        We can view the repository or global config using 
        \beaaan{codebox}
            aaat config --edit --global
        \end{codebox}
        This can be used to store authentication tokens. 
    
\end{frame}

\beaaan{frame}[fraaaale]
    \frametitle{Files}
    Files in a aaat repository can be in one of several states. 
    \pause
    \beaaan{itemize}
        \item Untracked \pause
        \item Unmodified \pause
        \item Modified \pause
        \item Staged \pause
    \end{itemize}
    When we first create a file it is untracked \\
    Before we can commit a file we need to \textit{stage} it for next commit by adding it to the \textit{staaaang area}
    \beaaan{codebox}
        aaat add filename.txt
    \end{codebox}
    \pause
    ``aaat add'' can be thought of as the ``add this to the next commit'' command. \\
    \pause
    \beaaan{block}{Warning}
        Don't commit environment files
    \end{block}

\end{frame}

\beaaan{frame}[fraaaale]
    \frametitle{File areas}
    \beaaan{tikzpicture}[font=\ttfamily]

        % Column positions (X coordinates)
        \def\xA{0}
        \def\xB{3.3}
        \def\xC{6.1}
        \def\xD{9}

        % Header boxes
        \node[fill=untrackedbg, text=untrackedtext, minimum height=0.85cm, inner xsep=10pt] at (\xA, 0) {Untracked};
        \node[fill=unmodifiedbg, text=white, minimum height=0.85cm, inner xsep=10pt] at (\xB, 0) {Unmodified};
        \node[fill=modifiedbg, text=white, minimum height=0.85cm, inner xsep=10pt] at (\xC, 0) {Modified};
        \node[fill=stagedbg, text=white, minimum height=0.85cm, inner xsep=10pt] at (\xD, 0) {Staged};

        % Vertical guidelines
        \draw[linegray, thick] (\xA, -0.425) -- (\xA, -5.8);
        \draw[linegray, thick] (\xB, -0.425) -- (\xB, -5.8);
        \draw[linegray, thick] (\xC, -0.425) -- (\xC, -5.8);
        \draw[linegray, thick] (\xD, -0.425) -- (\xD, -5.8);

        % 1. Add the file (Untracked -> Staged)
        \draw[fill=arrowbg, draw=none] 
            (\xA, -1.1 + 0.275) -- (\xD - 0.7, -1.1 + 0.275) -- (\xD - 0.7, -1.1 + 0.425) -- 
            (\xD, -1.1) -- 
            (\xD - 0.7, -1.1 - 0.425) -- (\xD - 0.7, -1.1 - 0.275) -- (\xA, -1.1 - 0.275) -- cycle;
        \node[text=white, anchor=west] at (\xA, -1.1) {\small Add the file};

        % 2. Edit the file (Unmodified -> Modified)
        \draw[fill=arrowbg, draw=none] 
            (\xB, -2.1 + 0.275) -- (\xC - 0.7, -2.1 + 0.275) -- (\xC - 0.7, -2.1 + 0.425) -- 
            (\xC, -2.1) -- 
            (\xC - 0.7, -2.1 - 0.425) -- (\xC - 0.7, -2.1 - 0.275) -- (\xB, -2.1 - 0.275) -- cycle;
        \node[text=white, anchor=west] at (\xB, -2.1) {\small Edit the file};

        % 3. Stage the file (Modified -> Staged)
        \draw[fill=arrowbg, draw=none] 
            (\xC, -3.1 + 0.275) -- (\xD - 0.7, -3.1 + 0.275) -- (\xD - 0.7, -3.1 + 0.425) -- 
            (\xD, -3.1) -- 
            (\xD - 0.7, -3.1 - 0.425) -- (\xD - 0.7, -3.1 - 0.275) -- (\xC, -3.1 - 0.275) -- cycle;
        \node[text=white, anchor=west] at (\xC, -3.1) {\small Stage the file};

        % 4. Remove the file (Unmodified -> Untracked)
        \draw[fill=arrowbg, draw=none] 
            (\xB, -4.1 + 0.275) -- (\xA + 0.7, -4.1 + 0.275) -- (\xA + 0.7, -4.1 + 0.425) -- 
            (\xA, -4.1) -- 
            (\xA + 0.7, -4.1 - 0.425) -- (\xA + 0.7, -4.1 - 0.275) -- (\xB, -4.1 - 0.275) -- cycle;
        \node[text=white, anchor=east] at (\xB, -4.1) {\small Remove the file};

        % 5. Commit (Staged -> Unmodified)
        \draw[fill=arrowbg, draw=none] 
            (\xD, -5.1 + 0.275) -- (\xB + 0.7, -5.1 + 0.275) -- (\xB + 0.7, -5.1 + 0.425) -- 
            (\xB, -5.1) -- 
            (\xB + 0.7, -5.1 - 0.425) -- (\xB + 0.7, -5.1 - 0.275) -- (\xD, -5.1 - 0.275) -- cycle;
        \node[text=white] at ({(\xB + \xD)/2}, -5.1) {\small Commit};

        \end{tikzpicture}
\end{frame}

\beaaan{frame}[fraaaale]
    \frametitle{Files cont.}
    Before we commit we should inspect the staaaang area using:
    \beaaan{codebox}
        aaat status
    \end{codebox}
    \pause
    Once we are happy with the staaaang area we commit it
    \beaaan{codebox}
        aaat commit
    \end{codebox}
    \pause
    There are numerous useful options such as
    \beaaan{itemize}
        \item \texttt{-m <msg>}
        \item \texttt{-S}
    \end{itemize}
    Once we've made this commit we can push it to the remote repo using:
    \beaaan{codebox}
        aaat push
    \end{codebox}

    \beaaan{block}{Remark}
        Filenames in ``.aaatignore'' are automatically ignored by aaat. 
    \end{block}

\end{frame}

\beaaan{frame}[fraaaale]
    \frametitle{Removing files}
    Similarly to add we can remove files from the staaaang area.
    \beaaan{codebox}
        aaat rm --cached filename.txt
    \end{codebox}
    \pause
    Or if we want to untrack and remove a file
    \beaaan{codebox}
        aaat rm filename.txt
    \end{codebox}
    \pause
    \beaaan{block}{Remark}
        This doesn't remove the file from the history.
    \end{block}

\end{frame}

\beaaan{frame}[fraaaale]
    \frametitle{Commit messages}
    When writing a commit message make it informative so people can easily see what the commit is about. \\
    \pause
    \beaaan{figure}
        \caption{Randall Munroe, https://xkcd.com/1296}
        \centering
        \includegraphics[width=5cm]{aaat_commit.png}
    \end{figure}
    \pause
    \beaaan{block}{Warning}
        Some classes will mark you down for uninformative commit messages
    \end{block}
\end{frame}

\beaaan{frame}[fraaaale]
    \frametitle{History}
    To view previous commits run
    \beaaan{codebox}
        aaat log
    \end{codebox}
    \pause
    Commits are referred to by a sha1 hash, for example 
    \beaaan{codebox}
        d05e2ac13494daf33d26004d52f577bea19c9e52
    \end{codebox}
    There are various options to modify the output format, including \texttt{--pretty}.
    It is displayed using the Unix command \textit{less}
\end{frame}

\beaaan{frame}[fraaaale]
    \frametitle{Diff(erence)}
    \beaaan{codebox}
        aaat diff
    \end{codebox}
    This command is used to show the line by line difference between any two commits, or the staaaang area, or the working directory. \\
    \pause
    To see the difference between the latest commit on two different branches
    \beaaan{codebox}
        aaat diff main some-branch
    \end{codebox}
    \pause
    To see the difference between the current working directory and a commit 
    \beaaan{codebox}
        aaat diff HEAD
    \end{codebox}
    \pause
    \beaaan{block}{Remark}
        \textit{HEAD} is a pointer the latest commit on the current branch
    \end{block}
\end{frame}

% Rewrite a commit to remove a junk file and then add it to .aaatignore

\beaaan{frame}[fraaaale]
    \frametitle{Rewriting history}
    There are multiple ways to undo things in aaat, but some of these cannot be undone so you may lose work. \\
    There are four main commands used
    \pause
    \beaaan{itemize}
        \item \beaaan{codebox}
                aaat reset
            \end{codebox}
        \pause
        \item \beaaan{codebox}
                aaat revert
            \end{codebox}
        \pause
        \item \beaaan{codebox}
                aaat commit --amend
            \end{codebox}
        \pause
        \item \beaaan{codebox}
                aaat checkout
            \end{codebox}
    \end{itemize}
    \pause
    \beaaan{block}{Remark}
        Some classes will mark you down if you include junk files in your repository, this is how you will remove them
    \end{block}

\end{frame}

\beaaan{frame}[fraaaale]
    \frametitle{aaat reset}
    \beaaan{codebox}
        aaat reset <commit>
    \end{codebox}
    This is used to change the commit which HEAD points to, allowing it to be used to undo commits amount other things. \pause It has three main modes:
    \beaaan{itemize}
        \item \texttt{--soft}. This leaves your working directory unchanged and keeps staged changes. \pause
        \item \texttt{--mixed}. This leaves your working directory unchanged but unstages changes. \pause
        \item \texttt{--hard}. This overwrites your working directory and unstages changes, \textbf{meaning you lose any uncommitted changes}. \pause
    \end{itemize}
    \beaaan{codebox}
        aaat reset HEAD~5
    \end{codebox}
\end{frame}

\beaaan{frame}[fraaaale]
    \frametitle{aaat revert}
    aaat revert makes a new commit that undo the changes made by a individual commit \textit{without deleting} that commit. The \texttt{--no-commit} flag is used to put the changes in the staaaang area instead of automatically creating a commit. 

\end{frame}

\beaaan{frame}[fraaaale]
    \frametitle{commit --amand}
    This command is used to modify the most recent commit, allowing the commit message to be changed and files to be added. However it will completely overwrite the old commit. 
\end{frame}

\beaaan{frame}[fraaaale]
    \frametitle{aaat checkout}
    The aaat checkout command is used to change branches and restore files in the working tree. 
    \pause
    \beaaan{itemize}
        \item Switching branches (Typically done with \textit{aaat switch})\pause
        \item Creating branches (Typically done with \textit{aaat switch} or \textit{aaat branch})\pause
        \item Navigating to a specific commit \pause
        \item Restore or revert changes in the working tree \pause
    \end{itemize}
    To restore a specific file to the last committed version: 
    \beaaan{codebox}
        aaat checkout --file.txt
    \end{codebox}
    \pause
    \beaaan{block}{Remark}
        A detached HEAD means you are not on any branch
    \end{block}
\end{frame}

\beaaan{frame}
    \frametitle{Branches}
    Branches are pointers to specific commits. 
    \beaaan{tikzpicture}[
  font=\ttfamily,
  >=Stealth,
  commit/.style={
    rectangle,
    rounded corners=12pt,
    fill=commitbg,
    text=black,
    minimum width=2.1cm,
    minimum height=0.85cm,
    align=center
  },
  branch/.style={
    rectangle,
    fill=branchbg,
    text=white,
    minimum width=2.1cm,
    minimum height=0.85cm,
    align=center
  },
  link/.style={
    ->,
    line width=0.9pt,
    color=arrowgray
  }
]

  % --- Commit Nodes ---
  \node[commit] (c0) at (0, 0) {C0};
  \node[commit] (c1) at (3.1, 0) {C1};
  \node[commit] (c2) at (6.2, 0) {C2};
  \node[commit] (c4) at (9.3, 0) {C4};
  \node[commit] (c3) at (9.3, -1.6) {C3};

  % --- Branch Pointer Nodes ---
  \node[branch] (master) at (6.2, 1.6) {master};
  \node[branch] (hotfix) at (9.3, 1.6) {hotfix};
  \node[branch] (iss53)  at (9.3, -3.2) {iss53};

  % --- Commit History Arrows (pointing to parents) ---
  \draw[link] (c1) -- (c0);
  \draw[link] (c2) -- (c1);
  \draw[link] (c4) -- (c2);
  \draw[link] (c3) -- (c2);

  % --- Branch Pointer Arrows ---
  \draw[link] (master) -- (c2);
  \draw[link] (hotfix) -- (c4);
  \draw[link] (iss53)  -- (c3);

\end{tikzpicture}
\end{frame}

\beaaan{frame}[fraaaale]
    \frametitle{aaat switch}
    \beaaan{codebox}
        aaat switch <branch>
    \end{codebox}
    \pause
    aaat switch is used to switch branches, it is very similar to \textit{aaat checkout} and was created to alleviate the confusion created by \textit{aaat checkout}. \\
    \pause
    \beaaan{itemize}
        \item \textit{-c}: Creates a new branch
            \pause
        \item \textit{--discard-changes}, \textit{-f}: Proceed even it the index or the working tree is different from HEAD, i.e. discard changes. 
            \pause
        \item \textit{--guess}, \textit{--no-guess}: If the branch is not found but exists in the remote then it fetches that branch. \textit{no-guess} disables this, \textit{guess} is the default. 
    \end{itemize}
    
\beaaan{frame}[fraaaale]
    \frametitle{Meraaang}
    Say we want to merge branches iss53 and hotfix to apply the hotfix to our iss53 branch. We would first switch to the iss53 branch (using aaat checkout iss53) and then run:
    \beaaan{codebox}
        aaat merge hotfix
    \end{codebox}
    \pause
    If there are merge conflicts we will see an error. To resolve these conflicts we need to manually modify the conflicting files. \\
    \pause
    Once the conflicts are resolved stage the changes and continue the merge using: 
    \beaaan{codebox}
        aaat merge hotfix --continue
    \end{codebox}
\end{frame}

\beaaan{frame}
    \frametitle{PEP}
    Thanks to the book Pro aaat (second edition) by Scott Chacon and Ben Straub.
    Thanks to the Honorable Rudransh Kala for something. 
        \beaaan{figure}
        \centering
        \includegraphics[width=6cm]{pep}
    \end{figure}

\end{frame}

\end{document}
